% Handbook 3.1, 10, 13. Part 1 owns data definitions.
\section{Problem and Data Description}
\label{sec:data}

\subsection{Problem formulation and research questions}
Given a grayscale image $x\in\mathbb{R}^{1\times28\times28}$, predict a label
$y\in\{0,\ldots,9\}$ identifying its clothing category. The prediction unit is one
image. A model $f_\theta$ produces ten logits, with prediction
$\hat y=\operatorname*{arg\,max}_{k} f_\theta(x)_k$.
\reporttodo{Explain a realistic motivation and the main difficulties, including
visually similar garment classes and the loss of spatial structure when flattening.
State research questions about predictive quality, representation/inductive bias,
and quality--compute trade-offs. Avoid promising a winning architecture in advance.}

\subsection{Dataset description and scope}
Fashion-MNIST is the mandatory primary evaluation dataset. MNIST is reserved for
development/debugging, and CIFAR-10 is an optional extension. Extension results
must be reported separately from the main Fashion-MNIST comparison.
\begin{table}[htbp]
\centering\small
\caption{Dataset documentation to verify against the dataset source.}
\label{tab:dataset}
\begin{tabularx}{\linewidth}{lY}
\toprule
Field & Description\\\midrule
Task / labels & Ten-class clothing image classification\\
Input / format & $28\times28$ grayscale; document original storage and loaded tensor dtype\\
Official sample counts & 60,000 training images; 10,000 test images\\
Source / version / license & \pending; add verified source and bibliographic entry\\
Collection / annotations & \pending; describe known collection and labeling process\\
Storage footprint & \pending; distinguish compressed download and local extracted size\\
Bias / limitations & \pending; discuss representativeness and real-world generalization\\
\bottomrule
\end{tabularx}
\end{table}
\reporttodo{List all ten label names in ID order. Cite the official dataset source.
Explain why MNIST and CIFAR-10 do not replace Fashion-MNIST in the core study.}

\subsection{Exploratory data analysis}
\label{sec:eda}
At minimum, report class distribution, input dimensions, imbalance analysis,
and representative samples. State the split used to compute each statistic.
\reportfigure{Images/eda/class_distribution.png}{Fashion-MNIST class distribution.
Source: the repository EDA export; identify the split and counts in the discussion.}{fig:class-distribution}
\reportfigure{Images/eda/representative_samples.png}{Representative raw Fashion-MNIST
samples. Source: the repository EDA export.}{fig:raw-samples}
\reporttodo{Quantify class balance with counts and an imbalance ratio; interpret
within-class variation and between-class similarities. Add other exported EDA
figures only when they support a specific decision or hypothesis. Check every
number against the EDA summary, rather than estimating it from a chart.}

\subsection{Data splitting and leakage prevention}
\label{sec:splits}
\begin{table}[htbp]
\centering
\caption{Split protocol. Counts below are the repository draft plan; verify the
actual loader and split indices before treating them as final.}
\label{tab:splits}
\begin{tabularx}{\linewidth}{l r Y}
\toprule
Split & Images & Role\\\midrule
Training & 54,000 & Model parameter fitting\\
Validation & 6,000 & Hyperparameter and checkpoint selection\\
Test & 10,000 & Final evaluation under a fixed protocol\\
\bottomrule
\end{tabularx}
\end{table}
\reporttodo{Document official versus custom partitioning, seed 42, split unit,
actual stratification policy and class counts for every split. Confirm that all
main models use identical indices. Explain leakage prevention, handling of
duplicates if checked, and subset rules if any. Distinguish the official
train/test distribution from the custom train/validation distribution.}

\subsection{Preprocessing, augmentation, and example batch}
\label{sec:preprocessing}
\reporttodo{Describe cleaning, invalid data/label handling, tensor conversion,
normalization constants and their rationale, and augmentation parameters.
Distinguish training transforms from deterministic validation/test transforms.
State whether resizing is needed and how labels are encoded.}
\begin{table}[htbp]
\centering\small
\caption{Expected representation shapes; verify against an actual batch.}
\label{tab:representations}
\begin{tabularx}{\linewidth}{l l Y}
\toprule
Representation & Shape & Models / interpretation\\\midrule
Image tensor & $B\times1\times28\times28$ & CNN; original grayscale geometry\\
Flattened pixels & $B\times784$ & Linear and MLP\\
Row sequence & $B\times28\times28$ & Recurrent models; 28 steps with 28 features\\
Tokens & $B\times T\times d$ & Transformer; define $T$, projection, and positions\\
\bottomrule
\end{tabularx}
\end{table}
\reportfigure{Images/eda/preprocessed_batch.png}{An actual training batch after
preprocessing; annotate tensor shape, dtype, label IDs, and intensity range.
This figure must be supplied from the loader.}{fig:processed-batch}
